\section{Exact finite-size PEPS representations}
\label{sec:peps_exact_tree}

The finite-size argument in~\cite[Section~8]{OpenAI2026PolynomialPEPS}
(\texttt{07-assembly.tex}, lines 203--214)
uses a spanning tree to represent every vector exactly. Let $G=(V,E)$ be a
finite connected graph, let $q\ge 1$, and put
$\Sigma=\{0,\ldots,q-1\}^{V}$ and $N=|\Sigma|=q^{|V|}$.

\begin{definition}[Configuration-copy tensors]
    \label{def:peps_configuration_copy}
    \lean{TNLean.PEPS.ExactTreeRepresentation.configurationEquiv,TNLean.PEPS.ExactTreeRepresentation.bondDim,TNLean.PEPS.ExactTreeRepresentation.decode,TNLean.PEPS.ExactTreeRepresentation.encode,TNLean.PEPS.ExactTreeRepresentation.Ports,TNLean.PEPS.ExactTreeRepresentation.Ports.ofConnected,TNLean.PEPS.ExactTreeRepresentation.localLabel,TNLean.PEPS.ExactTreeRepresentation.IsConsistent,TNLean.PEPS.ExactTreeRepresentation.tensor}
    \leanok
    Suppose $|V|\ge2$. Choose a connected spanning subgraph $T\subseteq G$, a
    root $r\in V$, and, at every vertex $v$, an incident edge $p(v)$ of $T$.
    For $\psi\colon\Sigma\to\C$, assign bond dimensions
    $D_e=N$ when $e\in T$ and $D_e=1$ otherwise. Decode each active edge index
    as a configuration $\tau_e\in\Sigma$, and define
    \begin{align}
        A_v((\tau_e)_{e\ni v},i)
         & = a_v(\tau_{p(v)})\,
        \delta_{i,\tau_{p(v)}(v)}
        \prod_{\substack{e\ni v     \\e\in T}}
        \delta_{\tau_e,\tau_{p(v)}},
         & a_v(\tau)            & =
        \begin{cases}\psi(\tau),&v=r,\\1,&v\ne r.\end{cases}
        \label{eq:peps_configuration_copy}
    \end{align}
    The one-dimensional indices on $E\setminus T$ do not enter the components.
\end{definition}

% Three-vertex path specialization of (eq:peps_configuration_copy).
% Ink-to-index: the two joined virtual wires each carry one Sigma index;
% the three open physical legs are the three local basis indices.
% Boundary signature: no open virtual indices and three open physical indices.
% No conjugation, trace, or normalization factor is present. Only the root
% tensor contains psi; all other tensors have entries zero or one.
\begin{center}
    \tenkzkernel
    \begin{tenkz}[rows={wire}, cols=3, bonds=none]
        \tn[at=(1,1), name=exactRoot, label pos=270,
            ports={0:virtual, 90:physical:$i_r$}]{A_r}
        \tn[at=(1,2), name=exactMiddle, label pos=270,
            ports={180:virtual, 0:virtual, 90:physical:$i_2$}]{A_2}
        \tn[at=(1,3), name=exactLast, label pos=270,
            ports={180:virtual, 90:physical:$i_3$}]{A_3}
        \tnwire[name=exactFirstBond]{exactRoot.0}{exactMiddle.180}
        \tnmark[form=label, label pos=n]{on exactFirstBond 0.5}{$\tau_{r2}$}
        \tnwire[name=exactSecondBond]{exactMiddle.0}{exactLast.180}
        \tnmark[form=label, label pos=n]{on exactSecondBond 0.5}{$\tau_{23}$}
    \end{tenkz}
\end{center}
For a three-vertex path, $\tau_{r2},\tau_{23}\in\Sigma$ are the two
contracted configuration indices; the open legs carry $i_r,i_2,i_3$.
Only $A_r$ contains the coefficient $\psi$. The same definition applies
to a vertex of any degree.

\begin{theorem}[Uniqueness of the contributing configuration]
    \label{thm:peps_copy_unique}
    \lean{TNLean.PEPS.ExactTreeRepresentation.decode_encode,TNLean.PEPS.ExactTreeRepresentation.decode_injective,TNLean.PEPS.ExactTreeRepresentation.localLabel_encode,TNLean.PEPS.ExactTreeRepresentation.isConsistent_encode,TNLean.PEPS.ExactTreeRepresentation.tensor_component_encode,TNLean.PEPS.ExactTreeRepresentation.isConsistent_of_component_ne_zero,TNLean.PEPS.ExactTreeRepresentation.localLabel_eq_of_adj,TNLean.PEPS.ExactTreeRepresentation.localLabel_eq_of_consistent,TNLean.PEPS.ExactTreeRepresentation.eq_encode_of_consistent,TNLean.PEPS.ExactTreeRepresentation.tensor_product_eq_zero}
    \leanok
    \uses{def:peps_configuration_copy}
    Fix $\sigma\in\Sigma$. If all factors in
    $\prod_{v\in V}A_v((\tau_e)_{e\ni v},\sigma(v))$ are nonzero, then
    $\tau_e=\sigma$ for every $e\in T$. Conversely, this virtual assignment
    satisfies all equality constraints.
\end{theorem}
\begin{proof}
    \leanok
    For an active edge $e=\{u,v\}$, the local constraints give
    $\tau_{p(u)}=\tau_e=\tau_{p(v)}$. Connectivity propagates these equalities
    along a walk between any two vertices. Hence there is one common
    configuration $\tau$. The physical constraint at each $v$ gives
    $\tau(v)=\sigma(v)$, so $\tau=\sigma$. The inactive edges have a unique
    index because their dimension is one. Substituting $\tau_e=\sigma$
    verifies the converse directly.
\end{proof}

\begin{theorem}[Exact contraction on a spanning tree]
    \label{thm:peps_exact_tree}
    \lean{TNLean.PEPS.ExactTreeRepresentation.stateCoeff_tensor,TNLean.PEPS.ExactTreeRepresentation.stateCoeff_tensor_eq,TNLean.PEPS.ExactTreeRepresentation.stateCoeff_tensor_ne_zero,TNLean.PEPS.ExactTreeRepresentation.bondDim_pos,TNLean.PEPS.ExactTreeRepresentation.bondDim_le,TNLean.PEPS.ExactTreeRepresentation.exists_exact_tree_tensor}
    \leanok
    \uses{def:peps_configuration_copy}
    For every finite connected graph with $|V|\ge2$ and every
    $\psi\colon\Sigma\to\C$, there are a spanning tree $T\subseteq G$ and
    tensors with dimensions $D_e=N$ on $T$ and $D_e=1$ off $T$ such that
    \begin{align}
        c_A(\sigma)
         & =\sum_{\eta}\prod_{v\in V} A_v(\eta_v,\sigma(v))
        =\psi(\sigma).
        \label{eq:peps_exact_tree_contraction}
    \end{align}
    Every bond is positive and at most $q^{|V|}$. If $\psi\ne0$, then $c_A\ne0$.
\end{theorem}
\begin{proof}
    \leanok
    \uses{thm:peps_copy_unique}
    Choose a spanning tree of $G$. Its connectedness and $|V|\ge2$ provide
    an incident tree edge at each vertex. By
    Theorem~\ref{thm:peps_copy_unique}, all terms of the contraction vanish
    except the assignment $\tau_e=\sigma$. Its root factor is $\psi(\sigma)$
    and every other factor is one, giving
    (\ref{eq:peps_exact_tree_contraction}). Since $q\ge1$, both possible bond
    dimensions are positive and at most $N$.
\end{proof}

\begin{theorem}[Singleton and empty vertex sets]
    \label{thm:peps_exact_degenerate}
    \lean{TNLean.PEPS.ExactTreeRepresentation.singletonTensor,TNLean.PEPS.ExactTreeRepresentation.stateCoeff_singletonTensor,TNLean.PEPS.ExactTreeRepresentation.exists_exact_tensor,TNLean.PEPS.ExactTreeRepresentation.stateCoeff_of_isEmpty}
    \leanok
    \uses{thm:peps_exact_tree}
    A connected graph with one vertex represents every coefficient function by
    its single physical tensor. Thus every finite connected graph admits the
    exact representation with positive bonds bounded by $q^{|V|}$.
    For an empty vertex set, however, the contraction is the scalar one for
    every tensor family.
\end{theorem}
\begin{proof}
    \leanok
    \uses{thm:peps_exact_tree}
    For one vertex $v$, set $A_v(i)=\psi(\sigma_i)$, where
    $\sigma_i(v)=i$. The graph has no edges, so the contraction is
    $A_v(\sigma(v))=\psi(\sigma)$. For no vertices there is one empty virtual
    assignment, and its empty product is one.
\end{proof}

\begin{theorem}[Connectivity of the open square grid]
    \label{thm:peps_square_connected}
    \lean{TNLean.PEPS.squareLatticeGraph_eq_boxProd,TNLean.PEPS.squareLatticeGraph_preconnected,TNLean.PEPS.squareLatticeGraph_connected,TNLean.PEPS.squareLatticeGraph_connected_iff,TNLean.PEPS.squareLatticeVertex_nontrivial}
    \leanok
    The open $w\times h$ nearest-neighbor grid is the Cartesian product of the
    path graphs on $w$ and $h$ vertices. It is connected exactly when $w,h>0$.
    If $L\ge2$, its $L\times L$ vertex set has at least two elements.
\end{theorem}
\begin{proof}
    \leanok
    The product adjacency condition says that one coordinate agrees and the
    other differs by one. Each coordinate path is preconnected, and
    preconnectedness is preserved by the Cartesian product. The vertex set is
    nonempty exactly when both side lengths are positive. For $L\ge2$, the
    vertices $(0,0)$ and $(1,0)$ are distinct.
\end{proof}

\begin{theorem}[Exact square-grid tensors and normalization]
    \label{thm:peps_exact_square}
    \lean{TNLean.PEPS.ExactTreeRepresentation.exists_exact_square_tensor,TNLean.PEPS.ExactTreeRepresentation.exists_exact_square_tree_tensor,TNLean.PEPS.ExactTreeRepresentation.exists_exact_unit_square_tensor}
    \leanok
    \uses{thm:peps_exact_degenerate,thm:peps_square_connected}
    For $q,L\ge1$, every vector $\psi$ on the open $L\times L$ square grid
    has an exact PEPS with $1\le D_e\le q^{L^2}$. For $L\ge2$, the dimensions
    can be chosen as $q^{L^2}$ on a spanning tree and one off that tree.
    If $\|\psi\|=1$, its contraction $\Phi$ satisfies
    $\Phi=\psi$, $\|\Phi\|=1$, and
    $\|\|\Phi\|^{-1}\Phi-1\cdot\psi\|=0$.
\end{theorem}
\begin{proof}
    \leanok
    \uses{thm:peps_exact_tree,thm:peps_exact_degenerate,thm:peps_square_connected}
    The grid has $L^2$ vertices. Apply the connected-graph result and identify
    equal coefficient functions with equal vectors in the physical Hilbert
    space. For a unit input the normalization scalar is one.
\end{proof}

\begin{theorem}[One prefactor for all small square sizes]
    \label{thm:peps_exact_small_sizes}
    \lean{TNLean.PEPS.ExactTreeRepresentation.exists_exact_square_tensor_bounded,TNLean.PEPS.ExactTreeRepresentation.exists_exact_square_tensor_polynomial,TNLean.PEPS.ExactTreeRepresentation.exists_exact_square_tensor_outer_power}
    \leanok
    \uses{thm:peps_exact_square}
    Fix $q\ge1$ and $L_0\in\N$. For every $0<L\le L_0$, every vector on the
    open $L\times L$ square grid has an exact PEPS with positive bonds at most
    $C_0=q^{L_0^2}$. This bound is independent of the vector and any Hamiltonian. For any $C\in\R$ and $r\ge0$, the enlarged prefactor $\max\{C,C_0\}$ gives $D_e\le\max\{C,C_0\}L^r$ throughout the same range. If $\chi>0$, then
    $C_\chi=\max\{C,C_0^{1/\chi}\}$ gives
    $D_e\le(C_\chi L^r)^\chi$ on this range as well.
\end{theorem}
\begin{proof}
    \leanok
    \uses{thm:peps_exact_square}
    Since $L^2\le L_0^2$ and $q\ge1$, the bound
    $D_e\le q^{L^2}\le q^{L_0^2}$ follows. In particular the same constant
    covers the manuscript's finite range $2\le L<L_0$. Since $L^r\ge1$, multiplying the enlarged nonnegative prefactor by $L^r$ preserves the bound. For the outer
    power, $C_\chi L^r\ge C_0^{1/\chi}$; raising both nonnegative sides to
    $\chi>0$ gives $(C_\chi L^r)^\chi\ge C_0$.
\end{proof}
